\chapter{Reticoli}\label{cap:reticoli}

Un \textbf{reticolo} è un insieme ordinato in cui due elementi qualsiasi hanno sempre un estremo inferiore e un estremo superiore.
Gli esempi principali li conosciamo già: le parti di un insieme con $\subseteq$, dove inf e sup sono intersezione e unione, e i naturali
con la divisibilità, dove sono MCD e mcm. In questo capitolo guardiamo i reticoli in due modi, come \textbf{insiemi ordinati} e come
\textbf{strutture algebriche} con due operazioni, e scopriamo che sono la stessa cosa. Poi studiamo sottoreticoli, complementi e
reticoli distributivi, che ci porteranno alle algebre di Boole.

% ──────────────────────────────────────────────────────────────
\section{Reticoli come insiemi ordinati}\label{sec:reticoli-ordinati}

\dfn[reticolo]{Reticolo}{
    Un insieme ordinato $(s, \rho)$, con $s \neq \varnothing$, si dice \textbf{reticolo} se per ogni $x, y \in s$ l'insieme $\{x, y\}$ ha estremo inferiore
    ed estremo superiore. In un reticolo sono quindi definite due operazioni interne, dette \textbf{operazioni reticolari}:
    \[ \wedge_\rho \colon (x, y) \in s \times s \mapsto \inf\{x, y\} \in s, \qquad\qquad \vee_\rho \colon (x, y) \in s \times s \mapsto \sup\{x, y\} \in s. \]
    $x \wedge y$ si legge «$x$ inf $y$» e $x \vee y$ si legge «$x$ sup $y$»; omettiamo il pedice $\rho$ quando non c'è ambiguità.
}

Per controllare che un insieme ordinato sia un reticolo basta guardare le coppie di elementi \textbf{non confrontabili}: se $x \rho y$,
l'insieme $\{x, y\}$ ha minimo $x$ e massimo $y$, che sono anche il suo inf e il suo sup (Proposizione~\ref{prop:minimo-inf}).

\begin{esempio}{Reticoli e non reticoli}
    \begin{center}
    \begin{tikzpicture}[hasse]
        \node[hv] (1) at (0,0) {$1$}; \node[hv] (2) at (-0.7,0.9) {$2$}; \node[hv] (3) at (0.7,0.9) {$3$};
        \node[hv] (4) at (-1.4,1.8) {$4$}; \node[hv] (6) at (0,1.8) {$6$}; \node[hv] (12) at (-0.7,2.7) {$12$};
        \draw[hlato] (1) -- (2) -- (4) -- (12); \draw[hlato] (1) -- (3) -- (6) -- (12); \draw[hlato] (2) -- (6);
        \node[font=\small, align=center] at (-0.35,-0.85) {divisori di $12$ con $\mid$:\\ è un reticolo};
        \begin{scope}[xshift=5.2cm]
            \node[hv evid] (a) at (-0.6,1.35) {$a$}; \node[hv evid] (b) at (0.6,1.35) {$b$};
            \node[font=\small, align=center] at (0,-0.85) {due elementi non confrontabili:\\ $\{a, b\}$ non ha maggioranti né minoranti};
        \end{scope}
        \begin{scope}[xshift=11.2cm]
            \node[hv evid] (a) at (-0.7,0.5) {$a$}; \node[hv evid] (b) at (0.7,0.5) {$b$};
            \node[hv evid2] (c) at (-0.7,2.2) {$c$}; \node[hv evid2] (d) at (0.7,2.2) {$d$};
            \draw[hlato] (a) -- (c); \draw[hlato] (a) -- (d); \draw[hlato] (b) -- (c); \draw[hlato] (b) -- (d);
            \node[font=\small, align=center] at (0,-0.85) {«farfalla»: $\{a, b\}$ ha i maggioranti $c$ e $d$,\\ ma nessuno dei due è il più piccolo};
        \end{scope}
    \end{tikzpicture}
    \end{center}
    Il secondo e il terzo insieme ordinato non sono reticoli: nel secondo $\{a, b\}$ non ha estremo superiore perché non ha proprio
    maggioranti; nel terzo i maggioranti di $\{a, b\}$ ci sono, ma $c$ e $d$ non sono confrontabili e quindi $\mathrm{Maggior}(\{a, b\})$ non ha minimo.
\end{esempio}

\dfn[limitato-completo]{Reticolo limitato, reticolo completo}{
    Un reticolo $(s, \rho)$ si dice \textbf{limitato} se ha minimo e massimo; si dice \textbf{completo} se \textbf{ogni} parte non vuota di $s$ (non solo
    quelle con due elementi) ha estremo inferiore ed estremo superiore.
}

\begin{concetto}
\begin{Proposizione}{Ogni reticolo completo è limitato}{completo-limitato}
    \begin{prerequisiti}
        \dip{Definizione}{def:limitato-completo}
        \dip{Definizione}{def:inf-sup}
        \dip{Definizione}{def:minimo-minimale}
    \end{prerequisiti}
    Se $(s, \rho)$ è un reticolo completo, allora $\max(s) = \sup(s)$ e $\min(s) = \inf(s)$; in particolare $(s, \rho)$ è limitato.
\end{Proposizione}
\begin{dimostrazione}
    Poiché $s \neq \varnothing$, esiste $M = \sup(s)$. $M$ è un maggiorante di $s$, cioè $x \rho M$ per ogni $x \in s$, e $M \in s$: quindi $M$ è il massimo di
    $s$. Allo stesso modo $\inf(s)$ è un minorante di $s$ che sta in $s$, cioè il minimo.
\end{dimostrazione}
\end{concetto}

\begin{spiegazione}
    Il sup di \textbf{tutto} $s$ sta sopra ogni elemento: è il «tetto» del reticolo. Il viceversa è falso: un reticolo limitato può non essere
    completo. Ad esempio $s = \{-1/n \mid n \geq 1\} \cup \{1/n \mid n \geq 1\}$, con l'ordine di $\mbQ$, è totalmente ordinato (quindi è un reticolo)
    e ha minimo $-1$ e massimo $1$. Però la parte $t = \{-1/n \mid n \geq 1\}$ non ha estremo superiore: i suoi maggioranti sono tutti i $1/n$, che
    non hanno minimo. L'estremo superiore «giusto» sarebbe $0$, che in $s$ manca.
    \begin{center}
    \begin{tikzpicture}
        \draw[black!45, line width=0.6pt] (-5.5,0) -- (5.5,0);
        \foreach \n in {1,...,10} {
            \fill[orange!85!black] ({-5/\n},0) circle (1.7pt);
            \fill[seagreen!70!black] ({5/\n},0) circle (1.7pt);
        }
        \draw[black!60, fill=white, line width=0.6pt] (0,0) circle (2pt);
        \node[below, font=\small] at (-5,-0.08) {$-1$}; \node[below, font=\small] at (-2.5,-0.08) {$-\frac{1}{2}$};
        \node[below, font=\small] at (-1.667,-0.08) {$-\frac{1}{3}$};
        \node[below, font=\small] at (5,-0.08) {$1$}; \node[below, font=\small] at (2.5,-0.08) {$\frac{1}{2}$}; \node[below, font=\small] at (1.667,-0.08) {$\frac{1}{3}$};
        \node[below, font=\small] at (0,-0.08) {$0 \notin s$};
        \node[font=\small, text=orange!80!black] at (-2.9,0.5) {$t = \{-1/n \mid n \geq 1\}$};
        \node[font=\small, text=seagreen!60!black] at (2.9,0.5) {maggioranti di $t$: i $1/n$, senza minimo};
    \end{tikzpicture}
    \end{center}
\end{spiegazione}

\begin{esempio}{Reticoli notevoli}
    \begin{itemize}
        \item $(\mcP(s), \subseteq)$ è un reticolo completo: per la Proposizione~\ref{prop:inf-parti}, $x \wedge y = x \cap y$ e $x \vee y = x \cup y$, e in generale
              $\inf(t) = \bigcap t$ e $\sup(t) = \bigcup t$. Il minimo è $\varnothing$, il massimo è $s$.
        \item $(\mbN, \mid)$ è un reticolo: $x \wedge y$ è il massimo comun divisore di $x$ e $y$ e $x \vee y$ il loro minimo comune multiplo (in $\mbN$ sono unici,
              Proposizione~\ref{prop:mcd-associati}). È limitato, con minimo $1$ (che divide tutti) e massimo $0$ (che è multiplo di tutti), ed è anche
              completo (Proposizione~\ref{prop:n-completo}).
        \item $(\mbN \setminus \{0\}, \mid)$ è un reticolo (MCD e mcm di numeri positivi sono positivi) con minimo $1$, ma \textbf{non ha massimo}: un massimo
              $M > 0$ dovrebbe essere multiplo di $2M$, che è più grande di $M$. Quindi non è limitato, e per la Proposizione~\ref{prop:completo-limitato} non è
              nemmeno completo.
        \item $(\mbQ, \leq)$, $(\mbZ, \leq)$ e in generale ogni insieme totalmente ordinato sono reticoli, con $x \wedge y = \min\{x, y\}$ e $x \vee y = \max\{x, y\}$.
              $(\mbZ, \leq)$ non è limitato, quindi non è completo: ad esempio la parte $\mbN$ non ha maggioranti in $\mbZ$.
    \end{itemize}
\end{esempio}

\warning{Attenzione allo zero}{
    Nel reticolo $(\mbN, \mid)$ il numero $0$ è il \textbf{massimo}, non il minimo: ogni naturale divide $0$. Togliendo lo $0$ si perde proprio il
    massimo, e $(\mbN \setminus \{0\}, \mid)$ non è più limitato.
}

% ──────────────────────────────────────────────────────────────
\Needspace{20\baselineskip}
\section{Dualità e minimali}\label{sec:dualita-reticoli}

Come per gli insiemi ordinati, capovolgendo un reticolo si ottiene ancora un reticolo; questo ci permette di dimostrare solo metà
degli enunciati.

\begin{concetto}
\begin{Proposizione}{Il duale di un reticolo è un reticolo}{reticolo-duale}
    \begin{prerequisiti}
        \dip{Definizione}{def:reticolo}
        \dip{Definizione}{def:relazione-duale}
        \dip{Proposizione}{prop:duale-proprieta}
    \end{prerequisiti}
    Se $(s, \rho)$ è un reticolo, anche $(s, \bar{\rho})$ è un reticolo, detto \textbf{reticolo duale}, e le due operazioni si scambiano:
    \[ x \wedge_{\bar{\rho}} y = x \vee_\rho y \qquad\qquad x \vee_{\bar{\rho}} y = x \wedge_\rho y. \]
\end{Proposizione}
\begin{dimostrazione}
    $(s, \bar{\rho})$ è un insieme ordinato (Proposizione~\ref{prop:duale-proprieta}). Passando da $\rho$ a $\bar{\rho}$ i minoranti diventano maggioranti:
    $\mathrm{Minor}_{\bar{\rho}}(\{x, y\}) = \mathrm{Maggior}_\rho(\{x, y\})$, e il massimo di questo insieme rispetto a $\bar{\rho}$ è il suo minimo rispetto a $\rho$.
    Quindi $\inf_{\bar{\rho}}\{x, y\}$ esiste ed è $\sup_\rho\{x, y\}$. Allo stesso modo $\sup_{\bar{\rho}}\{x, y\} = \inf_\rho\{x, y\}$.
\end{dimostrazione}
\end{concetto}

\begin{spiegazione}
    Capovolgendo il diagramma di Hasse ciò che stava sotto va sopra: il più grande dei minoranti diventa il più piccolo dei maggioranti. Nei divisori di $12$:
    \begin{center}
    \begin{tikzpicture}[hasse]
        \node[hv] (1) at (0,0) {$1$}; \node[hv evid2] (2) at (-0.7,0.9) {$2$}; \node[hv] (3) at (0.7,0.9) {$3$};
        \node[hv evid] (4) at (-1.4,1.8) {$4$}; \node[hv evid] (6) at (0,1.8) {$6$}; \node[hv evid2] (12) at (-0.7,2.7) {$12$};
        \draw[hlato] (1) -- (2) -- (4) -- (12); \draw[hlato] (1) -- (3) -- (6) -- (12); \draw[hlato] (2) -- (6);
        \node[font=\small, align=center] at (-0.35,-0.75) {con $\mid$: $4 \wedge 6 = 2$, $4 \vee 6 = 12$};
        \draw[-{Stealth}, line width=1pt, primary] (1.4,1.35) -- node[above, font=\small] {capovolgo} (3.0,1.35);
        \begin{scope}[xshift=5.0cm]
            \node[hv] (1) at (0,2.7) {$1$}; \node[hv evid2] (2) at (-0.7,1.8) {$2$}; \node[hv] (3) at (0.7,1.8) {$3$};
            \node[hv evid] (4) at (-1.4,0.9) {$4$}; \node[hv evid] (6) at (0,0.9) {$6$}; \node[hv evid2] (12) at (-0.7,0) {$12$};
            \draw[hlato] (1) -- (2) -- (4) -- (12); \draw[hlato] (1) -- (3) -- (6) -- (12); \draw[hlato] (2) -- (6);
            \node[font=\small, align=center] at (-0.35,-0.75) {con la duale: $4 \wedge_{\bar{\rho}} 6 = 12$, $4 \vee_{\bar{\rho}} 6 = 2$};
        \end{scope}
    \end{tikzpicture}
    \end{center}
\end{spiegazione}


\dfn[enunciato-duale-reticoli]{Enunciato duale per i reticoli}{
    Sia $e$ un enunciato sui reticoli. Il suo \textbf{duale} $\bar{e}$ si ottiene scambiando in $e$ $\rho$ con $\bar{\rho}$ e $\wedge$ con $\vee$ (e quindi minimo con massimo,
    minorante con maggiorante, inf con sup e così via).
}

\begin{concetto}
\begin{Teorema}{Principio di dualità per i reticoli}{dualita-reticoli}
    \begin{prerequisiti}
        \dip{Definizione}{def:enunciato-duale-reticoli}
        \dip{Proposizione}{prop:reticolo-duale}
    \end{prerequisiti}
    Se un enunciato $e$ vale per tutti i reticoli, allora anche il suo duale $\bar{e}$ vale per tutti i reticoli.
\end{Teorema}
\begin{dimostrazione}
    Sia $(s, \rho)$ un reticolo. Anche $(s, \bar{\rho})$ è un reticolo, quindi $e$ vale in $(s, \bar{\rho})$. Ma in $(s, \bar{\rho})$ l'ordine è $\bar{\rho}$ e gli inf e i sup
    sono i sup e gli inf di $(s, \rho)$ (Proposizione~\ref{prop:reticolo-duale}): dire che $e$ vale in $(s, \bar{\rho})$ è dire che $\bar{e}$ vale in $(s, \rho)$. Poiché
    $(s, \rho)$ è arbitrario, $\bar{e}$ vale per tutti i reticoli.
\end{dimostrazione}
\end{concetto}

\begin{spiegazione}
    Ogni volta che dimostriamo qualcosa su $\wedge$, otteniamo gratis l'enunciato «speculare» su $\vee$: sul diagramma di Hasse basta
    capovolgere il disegno. Alcune coppie di enunciati duali che incontreremo:
    \begin{center}
    \renewcommand{\arraystretch}{1.25}
    \begin{tabular}{l|l}
        \textbf{enunciato} & \textbf{enunciato duale} \\ \hline
        $(x \wedge y) \rho x$ & $x \rho (x \vee y)$ \\
        $x \wedge (y \wedge z) = (x \wedge y) \wedge z$ & $x \vee (y \vee z) = (x \vee y) \vee z$ \\
        $x \vee (x \wedge y) = x$ & $x \wedge (x \vee y) = x$ \\
        ogni elemento minimale è il minimo & ogni elemento massimale è il massimo \\
    \end{tabular}
    \end{center}
\end{spiegazione}

\begin{concetto}
\begin{Proposizione}{In un reticolo i minimali sono minimi}{minimale-minimo-reticolo}
    \begin{prerequisiti}
        \dip{Definizione}{def:minimo-minimale}
        \dip{Definizione}{def:reticolo}
        \dip{Teorema}{teo:dualita-reticoli}
    \end{prerequisiti}
    In un reticolo ogni elemento minimale è il minimo e ogni elemento massimale è il massimo.
\end{Proposizione}
\begin{dimostrazione}
    Per il principio di dualità basta la prima parte. Sia $m$ minimale e sia $x \in s$. L'elemento $m \wedge x$ è un minorante di $\{m, x\}$, quindi
    $(m \wedge x) \rho m$ e, per la minimalità di $m$, $m \wedge x = m$. Ma $m \wedge x$ sta anche sotto $x$: $m = m \wedge x \rho x$. Poiché $x$ era arbitrario, $m$ è il minimo.
\end{dimostrazione}
\end{concetto}

\begin{spiegazione}
    In un insieme ordinato qualunque ci possono essere più minimali (nella «farfalla» $a$ e $b$ lo sono entrambi). In un reticolo no: due
    minimali $m_1$ e $m_2$ hanno un inf, che sta sotto a entrambi; ma sotto un minimale c'è solo lui stesso, quindi $m_1 = m_1 \wedge m_2 = m_2$.
    \begin{center}
    \begin{tikzpicture}[hasse]
        \node[hv evid] (m1) at (-0.9,1.1) {$m_1$}; \node[hv evid] (m2) at (0.9,1.1) {$m_2$};
        \node[hv, dashed] (i) at (0,0) {$m_1 \wedge m_2$};
        \draw[hlato, dashed] (i) -- (m1); \draw[hlato, dashed] (i) -- (m2);
        \node[font=\small, align=left, anchor=west] at (1.9,0.55) {in un reticolo l'inf di $m_1$ e $m_2$ esiste e sta sotto entrambi:\\
            se $m_1$ e $m_2$ sono minimali deve essere $m_1 \wedge m_2 = m_1$ e $m_1 \wedge m_2 = m_2$.};
    \end{tikzpicture}
    \end{center}
    A differenza della Proposizione~\ref{prop:unico-minimale} non serve che il reticolo sia finito. Un reticolo può però non avere
    nessun minimale: in $(\mbZ, \leq)$ sotto ogni intero ce n'è un altro.
\end{spiegazione}

% ──────────────────────────────────────────────────────────────
\newpage
\section{Le parti finite di un reticolo}\label{sec:parti-finite-reticoli}

Per definizione, in un reticolo hanno inf e sup le parti con uno o due elementi. Vediamo che allora li hanno tutte le parti finite e non
vuote. Il lemma seguente vale in ogni insieme ordinato.

\begin{concetto}
\begin{Lemma}{Minoranti di un'unione}{minoranti-unione}
    \begin{prerequisiti}
        \dip{Definizione}{def:minoranti}
        \dip{Definizione}{def:inf-sup}
        \dip{Proposizione}{prop:impor}
        \dip{Proposizione}{prop:demorgan}
        \dip{Proposizione}{prop:distributivita}
        \dip{Proposizione}{prop:forall-and}
        \dip{Teorema}{teo:dualita-ordini}
    \end{prerequisiti}
    Sia $(s, \rho)$ un insieme ordinato e siano $a, b \subseteq s$.
    \begin{enumerate}
        \item $\mathrm{Minor}(a \cup b) = \mathrm{Minor}(a) \cap \mathrm{Minor}(b)$.
        \item Se $a$ ha estremo inferiore $m$, allora $\mathrm{Minor}(a) = \mathrm{Minor}(\{m\}) = \{x \in s \mid x \rho m\}$.
    \end{enumerate}
    Dualmente $\mathrm{Maggior}(a \cup b) = \mathrm{Maggior}(a) \cap \mathrm{Maggior}(b)$ e, se $M = \sup(a)$, $\mathrm{Maggior}(a) = \{x \in s \mid M \rho x\}$.
\end{Lemma}
\begin{dimostrazione}
    \begin{caso}{Punto 1}
        Siano $p$ = «$y \in a$», $q$ = «$y \in b$», $r$ = «$z \rho y$». La formula $(p \vee q) \rightarrow r$ equivale a $\neg(p \vee q) \vee r$, cioè a
        $(\neg p \wedge \neg q) \vee r$, cioè a $(\neg p \vee r) \wedge (\neg q \vee r)$, cioè a $(p \rightarrow r) \wedge (q \rightarrow r)$. Quindi
        \[ z \in \mathrm{Minor}(a \cup b) \leftrightarrow \forall y\,\big((p \vee q) \rightarrow r\big) \leftrightarrow \forall y\,(p \rightarrow r) \wedge \forall y\,(q \rightarrow r)
           \leftrightarrow z \in \mathrm{Minor}(a) \wedge z \in \mathrm{Minor}(b). \]
    \end{caso}
    \begin{caso}{Punto 2}
        Se $z \in \mathrm{Minor}(a)$, allora $z \rho m$, perché $m$ è il massimo dei minoranti. Viceversa, se $z \rho m$, per ogni $y \in a$ da $m \rho y$ segue
        $z \rho y$ (transitività), quindi $z \in \mathrm{Minor}(a)$. Infine $\mathrm{Minor}(\{m\}) = \{x \in s \mid x \rho m\}$ per definizione.
    \end{caso}
    Gli enunciati duali seguono dal principio di dualità (Teorema~\ref{teo:dualita-ordini}).
\end{dimostrazione}
\end{concetto}

\begin{spiegazione}
    Stare sotto a tutti gli elementi di $a \cup b$ vuol dire stare sotto a tutti quelli di $a$ \textbf{e} a tutti quelli di $b$. E i minoranti di un
    insieme sono esattamente gli elementi che stanno sotto il suo inf. Nei divisori di $12$, con $a = \{4\}$ e $b = \{6\}$:
    \begin{center}
    \renewcommand{\arraystretch}{1.2}
    \begin{tabular}{l|cccccc}
        $z$ & $1$ & $2$ & $3$ & $4$ & $6$ & $12$ \\ \hline
        $z \in \mathrm{Minor}(\{4\})$, cioè $z \mid 4$ & $\vero$ & $\vero$ & $\falso$ & $\vero$ & $\falso$ & $\falso$ \\
        $z \in \mathrm{Minor}(\{6\})$, cioè $z \mid 6$ & $\vero$ & $\vero$ & $\vero$ & $\falso$ & $\vero$ & $\falso$ \\
        $z \in \mathrm{Minor}(\{4, 6\})$ & $\vero$ & $\vero$ & $\falso$ & $\falso$ & $\falso$ & $\falso$ \\
        $z \in \mathrm{Minor}(\{2\})$, con $2 = \inf\{4, 6\}$ & $\vero$ & $\vero$ & $\falso$ & $\falso$ & $\falso$ & $\falso$ \\
    \end{tabular}
    \end{center}
    La terza riga è l'«e» delle prime due (punto 1) e coincide con la quarta (punto 2).
\end{spiegazione}

\begin{concetto}
\begin{Teorema}{Le parti finite di un reticolo hanno inf e sup}{parti-finite-reticolo}
    \begin{prerequisiti}
        \dip{Lemma}{lem:minoranti-unione}
        \dip{Proposizione}{prop:minimo-inf}
        \dip{Lemma}{lem:togli-elemento}
        \dip{Corollario}{cor:induzione-uso}
        \dip{Teorema}{teo:dualita-reticoli}
    \end{prerequisiti}
    Sia $(s, \rho)$ un reticolo. Ogni parte finita e non vuota di $s$ ha estremo inferiore ed estremo superiore. Inoltre, se una parte $u \subseteq s$
    ha estremo inferiore e $x \in s$, allora
    \[ \inf(u \cup \{x\}) = \inf(u) \wedge x \qquad\qquad\text{e, dualmente,}\qquad\qquad \sup(u \cup \{x\}) = \sup(u) \vee x. \]
\end{Teorema}
\begin{dimostrazione}
    Per il principio di dualità basta occuparsi dell'estremo inferiore.
    \begin{caso}{La formula}
        Sia $m = \inf(u)$. Per il Lemma~\ref{lem:minoranti-unione} (punto 1, poi punto 2, poi di nuovo punto 1)
        \[ \mathrm{Minor}(u \cup \{x\}) = \mathrm{Minor}(u) \cap \mathrm{Minor}(\{x\}) = \mathrm{Minor}(\{m\}) \cap \mathrm{Minor}(\{x\}) = \mathrm{Minor}(\{m, x\}), \]
        che ha massimo $m \wedge x$ perché $(s, \rho)$ è un reticolo. Quindi $\inf(u \cup \{x\}) = m \wedge x$.
    \end{caso}
    \begin{caso}{L'esistenza, per induzione su $n = |t| \geq 1$}
        \textbf{Base}: se $t = \{x\}$, $\inf(t) = x$ (Proposizione~\ref{prop:minimo-inf}). \textbf{Passo}: sia $|t| = n + 1 \geq 2$, sia $x \in t$ e sia $u = t \setminus \{x\}$,
        che ha $n \geq 1$ elementi (Lemma~\ref{lem:togli-elemento}). Per ipotesi induttiva $u$ ha estremo inferiore e, per la formula, anche $t = u \cup \{x\}$ ce l'ha.
    \end{caso}
\end{dimostrazione}
\end{concetto}

\begin{spiegazione}
    Il teorema dice come calcolare l'inf di una parte finita: si parte da un elemento e si aggiungono gli altri uno alla volta, proprio
    come in un ciclo \texttt{for} che aggiorna un accumulatore con $m \gets m \wedge x$. In $(\mbN, \mid)$, per $t = \{12, 18, 30\}$:
    \begin{passi}
        \item $\inf\{12\} = 12$ e $\sup\{12\} = 12$;
        \item aggiungo $18$: $\inf\{12, 18\} = 12 \wedge 18 = 6$ e $\sup\{12, 18\} = 12 \vee 18 = 36$;
        \item aggiungo $30$: $\inf\{12, 18, 30\} = 6 \wedge 30 = 6$ e $\sup\{12, 18, 30\} = 36 \vee 30 = 180$.
    \end{passi}
    Per le parti \textbf{infinite} il risultato è falso: $(\mbZ, \leq)$ è un reticolo, ma la sua parte $\mbN$ non ha estremo superiore.
\end{spiegazione}

\begin{concetto}
\begin{Corollario}{I reticoli finiti sono completi}{finito-completo}
    \begin{prerequisiti}
        \dip{Teorema}{teo:parti-finite-reticolo}
        \dip{Lemma}{lem:parti-finite}
        \dip{Proposizione}{prop:completo-limitato}
    \end{prerequisiti}
    Ogni reticolo finito è completo, e quindi limitato.
\end{Corollario}
\begin{dimostrazione}
    Se $s$ è finito, ogni sua parte non vuota è finita (Lemma~\ref{lem:parti-finite}) e quindi ha inf e sup (Teorema~\ref{teo:parti-finite-reticolo}): il reticolo è
    completo, e quindi limitato (Proposizione~\ref{prop:completo-limitato}).
\end{dimostrazione}
\end{concetto}

\begin{spiegazione}
    In un reticolo finito il minimo è l'inf di tutti gli elementi e il massimo è il loro sup: nei divisori di $12$ il minimo è
    $1 \wedge 2 \wedge 3 \wedge 4 \wedge 6 \wedge 12 = 1$ e il massimo è $1 \vee 2 \vee \dots \vee 12 = 12$. Un reticolo infinito, invece, può non essere limitato,
    come $(\mbZ, \leq)$.
\end{spiegazione}

\begin{concetto}
\begin{Proposizione}{$(\mbN, \mid)$ è un reticolo completo}{n-completo}
    \begin{prerequisiti}
        \dip{Proposizione}{prop:sottogruppi-z}
        \dip{Teorema}{teo:sottogruppi-generati}
        \dip{Teorema}{teo:intersezione-chiuse}
        \dip{Proposizione}{prop:prop-divisibilita}
        \dip{Definizione}{def:inf-sup}
    \end{prerequisiti}
    $(\mbN, \mid)$ è un reticolo completo, con minimo $1$ e massimo $0$. Per ogni parte non vuota $t \subseteq \mbN$:
    \begin{itemize}
        \item $\inf(t)$ è il numero $d \in \mbN$ tale che il sottogruppo di $(\mbZ, +)$ generato da $t$ è $d\mbZ$;
        \item $\sup(t)$ è il numero $\ell \in \mbN$ tale che $\bigcap_{x \in t} x\mbZ = \ell\mbZ$ (l'insieme dei multipli comuni degli elementi di $t$).
    \end{itemize}
\end{Proposizione}
\begin{dimostrazione}
    \facoltativa
    $1$ divide ogni naturale e ogni naturale divide $0$: sono il minimo e il massimo.
    \begin{caso}{Estremo inferiore}
        Il sottogruppo $\langle t \rangle$ è $d\mbZ$ per un $d \in \mbN$ (Proposizione~\ref{prop:sottogruppi-z}). Poiché $t \subseteq d\mbZ$, $d$ divide ogni elemento di $t$: è un
        minorante. Poiché $d \in \langle t \rangle$, $d$ è una somma di elementi di $t$ e di loro opposti (Teorema~\ref{teo:sottogruppi-generati}), quindi ogni divisore
        comune degli elementi di $t$ divide $d$ (Proposizione~\ref{prop:prop-divisibilita}): $d$ è il massimo dei minoranti.
    \end{caso}
    \begin{caso}{Estremo superiore}
        Ogni $x\mbZ$ è un sottogruppo di $(\mbZ, +)$, quindi lo è anche l'intersezione (Teorema~\ref{teo:intersezione-chiuse}), che allora è $\ell\mbZ$ con
        $\ell \in \mbN$. I maggioranti di $t$ in $(\mbN, \mid)$ sono i naturali multipli di tutti gli elementi di $t$, cioè i naturali di $\ell\mbZ$: $\ell$ è uno di
        loro e divide tutti gli altri, quindi è il minimo dei maggioranti.
    \end{caso}
\end{dimostrazione}
\end{concetto}

\begin{spiegazione}
    Il diagramma di Hasse di $(\mbN, \mid)$ ha $1$ in fondo, sopra i primi, poi i prodotti di due primi e così via; $0$ sta in cima a tutto,
    infinitamente in alto. Qualche estremo di parti \textbf{infinite}:
    \begin{center}
    \begin{tikzpicture}[hasse]
        \node[hv evid2] (1) at (0,0) {$1$};
        \node[hv] (2) at (-1.6,1) {$2$}; \node[hv] (3) at (0,1) {$3$}; \node[hv] (5) at (1.6,1) {$5$};
        \node[hv] (4) at (-2.6,2) {$4$}; \node[hv] (6) at (-1.3,2) {$6$}; \node[hv] (10) at (0,2) {$10$};
        \node[hv] (9) at (1.3,2) {$9$}; \node[hv] (15) at (2.6,2) {$15$};
        \draw[hlato] (1) -- (2); \draw[hlato] (1) -- (3); \draw[hlato] (1) -- (5);
        \draw[hlato] (2) -- (4); \draw[hlato] (2) -- (6); \draw[hlato] (2) -- (10); \draw[hlato] (3) -- (6); \draw[hlato] (3) -- (9);
        \draw[hlato] (3) -- (15); \draw[hlato] (5) -- (10); \draw[hlato] (5) -- (15);
        \node at (0,2.75) {$\vdots$};
        \node[hv evid2] (0) at (0,3.7) {$0$};
        \foreach \v in {4,6,10,9,15} \draw[hlato grigio] (\v) -- (0);
        \node[font=\small, align=left, anchor=west] at (3.6,1.85) {
            $\inf\{6, 10, 15\} = 1$: infatti $6 + 10 - 15 = 1$, quindi $\langle 6, 10, 15 \rangle = \mbZ$;\\[1mm]
            $\inf\{\text{pari positivi}\} = 2$ e $\inf\{4, 6, 8, 10, \dots\} = 2$;\\[1mm]
            $\sup\{2, 4, 8, 16, \dots\} = 0$: nessun numero non nullo\\ \hspace*{3mm} è multiplo di tutte le potenze di $2$;\\[1mm]
            in generale una parte infinita ha sup $0$: un numero\\ \hspace*{3mm} $\ell \neq 0$ ha solo un numero finito di divisori.};
    \end{tikzpicture}
    \end{center}
\end{spiegazione}

% ──────────────────────────────────────────────────────────────
\newpage
\section{Le operazioni reticolari}\label{sec:operazioni-reticolari}

Un reticolo $(s, \rho)$ dà luogo alla struttura algebrica $(s, \vee, \wedge)$, con due operazioni interne. Vediamone le proprietà.

\begin{concetto}
\begin{Teorema}{Proprietà delle operazioni reticolari}{proprieta-reticolari}
    \begin{prerequisiti}
        \dip{Teorema}{teo:parti-finite-reticolo}
        \dip{Proposizione}{prop:minimo-inf}
        \dip{Teorema}{teo:dualita-reticoli}
        \dip{Definizione}{def:commutativa-associativa}
    \end{prerequisiti}
    Sia $(s, \rho)$ un reticolo. Per ogni $x, y, z \in s$ valgono:
    \begin{enumerate}
        \item \textbf{commutatività}: $x \wedge y = y \wedge x$ e $x \vee y = y \vee x$;
        \item \textbf{associatività}: $x \wedge (y \wedge z) = (x \wedge y) \wedge z$ e $x \vee (y \vee z) = (x \vee y) \vee z$;
        \item \textbf{assorbimento}: $x \vee (x \wedge y) = x$ e $x \wedge (x \vee y) = x$;
        \item \textbf{idempotenza}: $x \wedge x = x$ e $x \vee x = x$.
    \end{enumerate}
\end{Teorema}
\begin{dimostrazione}
    In ogni punto le due uguaglianze sono l'una la duale dell'altra: per il principio di dualità basta dimostrare la prima.
    \begin{caso}{Punti 1 e 4}
        $\{x, y\} = \{y, x\}$, quindi i due insiemi hanno lo stesso inf. Inoltre $\{x, x\} = \{x\}$, quindi $x \wedge x = \inf\{x\} = x$ (Proposizione~\ref{prop:minimo-inf}).
    \end{caso}
    \begin{caso}{Punto 2}
        Per il Teorema~\ref{teo:parti-finite-reticolo}, $\inf\{x, y, z\} = \inf(\{y, z\} \cup \{x\}) = (y \wedge z) \wedge x$, e anche $\inf\{x, y, z\} = \inf(\{x, y\} \cup \{z\}) = (x \wedge y) \wedge z$.
        Per la commutatività $x \wedge (y \wedge z) = (y \wedge z) \wedge x = (x \wedge y) \wedge z$.
    \end{caso}
    \begin{caso}{Punto 3}
        $x \wedge y$ è un minorante di $\{x, y\}$, quindi $(x \wedge y) \rho x$: l'insieme $\{x, x \wedge y\}$ ha massimo $x$, e allora il suo estremo superiore è $x$
        (Proposizione~\ref{prop:minimo-inf}), cioè $x \vee (x \wedge y) = x$.
    \end{caso}
\end{dimostrazione}
\end{concetto}

\begin{spiegazione}
    Nell'insieme delle parti sono le proprietà di $\cap$ e $\cup$ già viste nella Proposizione~\ref{prop:prop-operazioni}. L'associatività dice che
    $x \wedge y \wedge z$ è sempre l'inf dei tre elementi, comunque si mettano le parentesi. L'assorbimento si vede bene con i diagrammi di Venn
    (il risultato è il bordo spesso):
    \begin{center}
    \begin{tikzpicture}
        \begin{scope}
            \begin{scope} \clip (0,0) circle (0.95); \fill[orange!40] (1.1,0) circle (0.95); \end{scope}
            \draw[black!55] (1.1,0) circle (0.95); \draw[orange!85!black, line width=1.6pt] (0,0) circle (0.95);
            \node at (-0.45,0) {$x$}; \node at (1.55,0) {$y$};
            \node[font=\small, align=center] at (0.55,-1.45) {$x \cup (x \cap y) = x$: unendo a $x$\\ un suo pezzo si ottiene $x$};
        \end{scope}
        \begin{scope}[xshift=6.5cm]
            \fill[primary!15] (0,0) circle (0.95); \fill[primary!15] (1.1,0) circle (0.95);
            \draw[black!55] (1.1,0) circle (0.95); \draw[orange!85!black, line width=1.6pt] (0,0) circle (0.95);
            \node at (-0.45,0) {$x$}; \node at (1.55,0) {$y$};
            \node[font=\small, align=center] at (0.55,-1.45) {$x \cap (x \cup y) = x$: intersecando $x$\\ con un insieme che lo contiene si ottiene $x$};
        \end{scope}
    \end{tikzpicture}
    \end{center}
    In $(\mbN, \mid)$: $12 \vee (12 \wedge 18) = \mathrm{mcm}(12, 6) = 12$ e $12 \wedge (12 \vee 18) = \mathrm{MCD}(12, 36) = 12$.
\end{spiegazione}

\begin{concetto}
\begin{Proposizione}{L'ordine si legge sulle operazioni}{ordine-operazioni}
    \begin{prerequisiti}
        \dip{Definizione}{def:reticolo}
        \dip{Proposizione}{prop:minimo-inf}
    \end{prerequisiti}
    Sia $(s, \rho)$ un reticolo e siano $x, y \in s$. Allora
    \[ x \rho y \quad\leftrightarrow\quad x \wedge y = x \quad\leftrightarrow\quad x \vee y = y. \]
\end{Proposizione}
\begin{dimostrazione}
    Se $x \rho y$, l'insieme $\{x, y\}$ ha minimo $x$ e massimo $y$, quindi $x \wedge y = x$ e $x \vee y = y$ (Proposizione~\ref{prop:minimo-inf}). Se $x \wedge y = x$,
    allora $x = \inf\{x, y\}$ è un minorante di $\{x, y\}$, quindi $x \rho y$. Se $x \vee y = y$, allora $y = \sup\{x, y\}$ è un maggiorante di $\{x, y\}$, quindi $x \rho y$.
\end{dimostrazione}
\end{concetto}

\begin{spiegazione}
    Chi conosce solo le operazioni può ricostruire l'ordine: $x$ sta sotto $y$ quando «$x$ inf $y$» non cambia $x$. Due casi familiari:
    \begin{itemize}
        \item in $(\mcP(s), \subseteq)$: $x \subseteq y \leftrightarrow x \cap y = x \leftrightarrow x \cup y = y$ (Proposizione~\ref{prop:caratt-inclusione});
        \item in $(\mbN, \mid)$: $4 \mid 12$, e infatti $\mathrm{MCD}(4, 12) = 4$ e $\mathrm{mcm}(4, 12) = 12$; invece $6 \nmid 4$, e infatti $\mathrm{MCD}(6, 4) = 2 \neq 6$.
    \end{itemize}
    Questa proposizione è il ponte tra i due punti di vista che confronteremo nella prossima sezione.
\end{spiegazione}

\begin{concetto}
\begin{Proposizione}{Elementi neutri delle operazioni reticolari}{neutri-reticolo}
    \begin{prerequisiti}
        \dip{Proposizione}{prop:ordine-operazioni}
        \dip{Definizione}{def:neutro}
        \dip{Definizione}{def:monoide}
    \end{prerequisiti}
    Sia $(s, \rho)$ un reticolo. Un elemento $e \in s$ è neutro rispetto a $\wedge$ se e solo se $e = \max(s)$, ed è neutro rispetto a $\vee$ se e solo se
    $e = \min(s)$. In particolare, se il reticolo è limitato, $(s, \wedge, \max(s))$ e $(s, \vee, \min(s))$ sono monoidi commutativi.
\end{Proposizione}
\begin{dimostrazione}
    Per la commutatività di $\wedge$ e la Proposizione~\ref{prop:ordine-operazioni}: $e$ è neutro per $\wedge$ $\leftrightarrow$ $\forall x \in s\,(x \wedge e = x)$ $\leftrightarrow$
    $\forall x \in s\,(x \rho e)$ $\leftrightarrow$ $e = \max(s)$. La parte su $\vee$ è duale. L'ultima affermazione segue allora dal Teorema~\ref{teo:proprieta-reticolari}.
\end{dimostrazione}
\end{concetto}

\begin{spiegazione}
    Il massimo non cambia nulla quando si prende un inf con lui, il minimo non cambia nulla quando si prende un sup:
    \begin{center}
    \renewcommand{\arraystretch}{1.25}
    \begin{tabular}{l|c|c}
        reticolo & neutro di $\wedge$ (il massimo) & neutro di $\vee$ (il minimo) \\ \hline
        $(\mcP(s), \subseteq)$ & $s$: \ $x \cap s = x$ & $\varnothing$: \ $x \cup \varnothing = x$ \\
        $(\mbN, \mid)$ & $0$: \ $\mathrm{MCD}(x, 0) = x$ & $1$: \ $\mathrm{mcm}(x, 1) = x$ \\
        $(\mbZ, \leq)$ & non esiste & non esiste \\
    \end{tabular}
    \end{center}
\end{spiegazione}

% ──────────────────────────────────────────────────────────────
\newpage
\section{Due punti di vista equivalenti}\label{sec:equivalenza-reticoli}

Le proprietà del Teorema~\ref{teo:proprieta-reticolari} bastano a ricostruire tutto il reticolo. Attenzione: nel lemma seguente $\vee$ e $\wedge$
sono due operazioni \textbf{qualsiasi} con quelle proprietà, non (ancora) gli inf e i sup di un ordine.

\begin{concetto}
\begin{Lemma}{Dalle operazioni all'ordine}{ordine-da-operazioni}
    \begin{prerequisiti}
        \dip{Definizione}{def:operazione-interna}
        \dip{Definizione}{def:tipi-relazioni}
        \dip{Definizione}{def:reticolo}
    \end{prerequisiti}
    Sia $s \neq \varnothing$ e siano $\vee$, $\wedge$ due operazioni interne su $s$ commutative, associative e per cui vale l'assorbimento:
    $x \vee (x \wedge y) = x$ e $x \wedge (x \vee y) = x$ per ogni $x, y \in s$. Definiamo su $s$ la relazione $x \rho y \defsse x = x \wedge y$. Allora:
    \begin{enumerate}
        \item $x \wedge x = x$ e $x \vee x = x$ per ogni $x$; inoltre $x = x \wedge y$ se e solo se $y = x \vee y$;
        \item $\rho$ è una relazione d'ordine su $s$;
        \item $(s, \rho)$ è un reticolo, con $\inf\{x, y\} = x \wedge y$ e $\sup\{x, y\} = x \vee y$.
    \end{enumerate}
\end{Lemma}
\begin{dimostrazione}
    \begin{caso}{Punto 1}
        Per il primo assorbimento $x \vee (x \wedge x) = x$; sostituendo nel secondo, $x \wedge x = x \wedge \big(x \vee (x \wedge x)\big) = x$. Allo stesso modo $x \vee x = x$.
        Se $x = x \wedge y$, allora $x \vee y = (x \wedge y) \vee y = y \vee (y \wedge x) = y$; se $y = x \vee y$, allora $x \wedge y = x \wedge (x \vee y) = x$.
    \end{caso}
    \begin{caso}{Punto 2}
        Riflessiva: $x = x \wedge x$. Antisimmetrica: da $x = x \wedge y$ e $y = y \wedge x$ segue $x = x \wedge y = y \wedge x = y$. Transitiva: da $x = x \wedge y$ e
        $y = y \wedge z$ segue $x \wedge z = (x \wedge y) \wedge z = x \wedge (y \wedge z) = x \wedge y = x$.
    \end{caso}
    \begin{caso}{Punto 3}
        $(x \wedge y) \wedge x = (x \wedge x) \wedge y = x \wedge y$ e $(x \wedge y) \wedge y = x \wedge (y \wedge y) = x \wedge y$: $x \wedge y$ è un minorante di $\{x, y\}$. Se $z$ è un
        minorante, cioè $z = z \wedge x$ e $z = z \wedge y$, allora $z \wedge (x \wedge y) = (z \wedge x) \wedge y = z \wedge y = z$, cioè $z \rho (x \wedge y)$. Quindi $x \wedge y = \inf\{x, y\}$.
        Infine, per il punto 1, $x \bar{\rho} y \leftrightarrow y = y \wedge x \leftrightarrow x = x \vee y$: $\bar{\rho}$ si ottiene da $\vee$ come $\rho$ da $\wedge$, e le ipotesi
        non cambiano scambiando $\wedge$ e $\vee$. Ripetendo il ragionamento, $x \vee y = \inf_{\bar{\rho}}\{x, y\} = \sup_\rho\{x, y\}$.
    \end{caso}
\end{dimostrazione}
\end{concetto}

\begin{spiegazione}
    Ogni proprietà dell'ordine viene da una precisa proprietà delle operazioni:
    \begin{center}
    \renewcommand{\arraystretch}{1.25}
    \begin{tabular}{l|l}
        \textbf{serve per} & \textbf{proprietà delle operazioni usata} \\ \hline
        riflessività di $\rho$ & idempotenza (che a sua volta viene dai due assorbimenti) \\
        antisimmetria di $\rho$ & commutatività \\
        transitività di $\rho$ & associatività \\
        $x \wedge y$ è l'inf & associatività, commutatività e idempotenza \\
        $x \vee y$ è il sup & punto 1 (cioè l'assorbimento) e lo stesso ragionamento con $\vee$ \\
    \end{tabular}
    \end{center}
    Il punto 1 dice anche che l'ordine si può leggere indifferentemente sulla tavola di $\wedge$ ($x \wedge y = x$) o su quella di $\vee$ ($x \vee y = y$).
\end{spiegazione}

\begin{concetto}
\begin{Teorema}{Reticoli come insiemi ordinati e come strutture algebriche}{reticoli-equivalenza}
    \begin{prerequisiti}
        \dip{Teorema}{teo:proprieta-reticolari}
        \dip{Proposizione}{prop:ordine-operazioni}
        \dip{Lemma}{lem:ordine-da-operazioni}
        \dip{Teorema}{teo:biettiva-inversa}
    \end{prerequisiti}
    Sia $s \neq \varnothing$. Sia $R(s)$ l'insieme delle relazioni d'ordine $\rho$ su $s$ tali che $(s, \rho)$ è un reticolo, e sia $B(s)$ l'insieme delle coppie
    $(\vee, \wedge)$ di operazioni interne su $s$ commutative, associative e per cui vale l'assorbimento. Allora la funzione
    \[ f \colon \rho \in R(s) \mapsto (\vee_\rho, \wedge_\rho) \in B(s) \]
    è biettiva, e la sua inversa associa a $(\vee, \wedge)$ l'ordine definito da $x \rho y \defsse x = x \wedge y$.
\end{Teorema}
\begin{dimostrazione}
    $f$ è ben definita per il Teorema~\ref{teo:proprieta-reticolari}. Sia $g$ la funzione che associa a $(\vee, \wedge) \in B(s)$ la relazione $x \rho y \defsse x = x \wedge y$:
    per il Lemma~\ref{lem:ordine-da-operazioni} è un elemento di $R(s)$, quindi $g \colon B(s) \to R(s)$.
    \begin{caso}{$f \circ g = \mathrm{id}_{B(s)}$}
        Per il Lemma~\ref{lem:ordine-da-operazioni}, inf e sup dell'ordine $g(\vee, \wedge)$ sono proprio $\wedge$ e $\vee$: $f(g(\vee, \wedge)) = (\vee, \wedge)$.
    \end{caso}
    \begin{caso}{$g \circ f = \mathrm{id}_{R(s)}$}
        Partendo da $\rho \in R(s)$, l'ordine $g(f(\rho))$ è «$x = x \wedge_\rho y$», che equivale a $x \rho y$ (Proposizione~\ref{prop:ordine-operazioni}): $g(f(\rho)) = \rho$.
    \end{caso}
    Quindi $g$ è l'inversa di $f$, e $f$ è biettiva (Teorema~\ref{teo:biettiva-inversa}).
\end{dimostrazione}
\end{concetto}

\begin{spiegazione}
    Lo stesso reticolo si può dare in due modi: con il diagramma di Hasse (l'ordine) o con le tavole delle due operazioni. Si passa da
    un modo all'altro nei due versi, e facendo andata e ritorno si torna al punto di partenza. Per i divisori di $6$
    (in arancione le caselle con $x \wedge y = x$, cioè le coppie con $x \mid y$):
    \begin{center}
    \begin{tikzpicture}[hasse]
        \node[hv] (1) at (0,0) {$1$}; \node[hv] (2) at (-0.8,1) {$2$}; \node[hv] (3) at (0.8,1) {$3$}; \node[hv] (6) at (0,2) {$6$};
        \draw[hlato] (1) -- (2) -- (6); \draw[hlato] (1) -- (3) -- (6);
        \node[font=\small] at (0,-0.7) {l'ordine $\mid$};
        \draw[-{Stealth}, line width=1pt, primary] (1.4,1.35) -- node[above, font=\small] {$x \wedge y := \inf\{x, y\}$} (4.6,1.35);
        \draw[{Stealth}-, line width=1pt, primary] (1.4,0.65) -- node[below, font=\small] {$x \rho y \defsse x = x \wedge y$} (4.6,0.65);
        \node[anchor=west] at (4.9,1.0) {\small
            \begin{tabular}{c|cccc}
                $\wedge$ & $1$ & $2$ & $3$ & $6$ \\ \hline
                $1$ & \cellcolor{orange!25}$1$ & \cellcolor{orange!25}$1$ & \cellcolor{orange!25}$1$ & \cellcolor{orange!25}$1$ \\
                $2$ & $1$ & \cellcolor{orange!25}$2$ & $1$ & \cellcolor{orange!25}$2$ \\
                $3$ & $1$ & $1$ & \cellcolor{orange!25}$3$ & \cellcolor{orange!25}$3$ \\
                $6$ & $1$ & $2$ & $3$ & \cellcolor{orange!25}$6$ \\
            \end{tabular}\quad
            \begin{tabular}{c|cccc}
                $\vee$ & $1$ & $2$ & $3$ & $6$ \\ \hline
                $1$ & $1$ & $2$ & $3$ & $6$ \\
                $2$ & $2$ & $2$ & $6$ & $6$ \\
                $3$ & $3$ & $6$ & $3$ & $6$ \\
                $6$ & $6$ & $6$ & $6$ & $6$ \\
            \end{tabular}};
        \node[font=\small] at (8.4,-0.7) {le operazioni: MCD e mcm};
    \end{tikzpicture}
    \end{center}
    Per questo d'ora in poi parleremo indifferentemente del reticolo $(s, \rho)$ o del reticolo $(s, \vee, \wedge)$.
\end{spiegazione}

\begin{concetto}
\begin{Teorema}{Isomorfismi di reticoli}{isomorfismi-reticoli}
    \begin{prerequisiti}
        \dip{Definizione}{def:crescente}
        \dip{Definizione}{def:omomorfismo}
        \dip{Proposizione}{prop:ordine-operazioni}
        \dip{Teorema}{teo:dualita-reticoli}
    \end{prerequisiti}
    Siano $(s, \rho)$ e $(s', \rho')$ reticoli e sia $f \colon s \to s'$. Allora $f$ è un isomorfismo tra gli insiemi ordinati $(s, \rho)$ e $(s', \rho')$ se e solo se è un
    isomorfismo tra le strutture algebriche $(s, \vee_\rho, \wedge_\rho)$ e $(s', \vee_{\rho'}, \wedge_{\rho'})$.
\end{Teorema}
\begin{dimostrazione}
    In entrambi i casi $f$ è biettiva; scriviamo $\wedge$, $\vee$ in $s$ e $\wedge'$, $\vee'$ in $s'$.
    \begin{caso}{$(\leftarrow)$ Un isomorfismo di strutture rispetta l'ordine}
        Per ogni $x, y \in s$: $x \rho y \leftrightarrow x = x \wedge y \leftrightarrow f(x) = f(x \wedge y) = f(x) \wedge' f(y) \leftrightarrow f(x) \rho' f(y)$, dove il secondo
        $\leftrightarrow$ usa l'iniettività di $f$ e gli altri la Proposizione~\ref{prop:ordine-operazioni}.
    \end{caso}
    \begin{caso}{$(\rightarrow)$ Un isomorfismo di insiemi ordinati rispetta le operazioni}
        Siano $x, y \in s$. Da $(x \wedge y) \rho x$ e $(x \wedge y) \rho y$ segue che $f(x \wedge y)$ è un minorante di $\{f(x), f(y)\}$. Sia $w$ un altro minorante; per la
        suriettività $w = f(z)$, e da $f(z) \rho' f(x)$, $f(z) \rho' f(y)$ segue $z \rho x$ e $z \rho y$, quindi $z \rho (x \wedge y)$ e $w = f(z) \rho' f(x \wedge y)$. Dunque
        $f(x \wedge y)$ è il massimo dei minoranti: $f(x \wedge y) = f(x) \wedge' f(y)$. Per $\vee$ il ragionamento è duale.
    \end{caso}
\end{dimostrazione}
\end{concetto}

\begin{spiegazione}
    Anche gli isomorfismi sono gli stessi nei due punti di vista. I divisori di $6$ e le parti di $\{a, b\}$ hanno lo stesso diagramma, e
    la biezione che li sovrappone trasforma MCD e mcm in intersezione e unione:
    \begin{center}
    \begin{tikzpicture}[hasse]
        \node[hv] (1) at (0,0) {$1$}; \node[hv evid] (2) at (-0.8,1) {$2$}; \node[hv evid] (3) at (0.8,1) {$3$}; \node[hv] (6) at (0,2) {$6$};
        \draw[hlato] (1) -- (2) -- (6); \draw[hlato] (1) -- (3) -- (6);
        \draw[-{Stealth}, line width=1pt, primary] (1.4,1) -- node[above, font=\small] {$f$} (2.8,1);
        \begin{scope}[xshift=4.4cm]
            \node[hv] (v) at (0,0) {$\varnothing$}; \node[hv evid] (a) at (-0.9,1) {$\{a\}$}; \node[hv evid] (b) at (0.9,1) {$\{b\}$};
            \node[hv] (ab) at (0,2) {$\{a, b\}$};
            \draw[hlato] (v) -- (a) -- (ab); \draw[hlato] (v) -- (b) -- (ab);
        \end{scope}
        \node[font=\small, align=left, anchor=west] at (6.2,1) {$f(1) = \varnothing$, \ $f(2) = \{a\}$, \ $f(3) = \{b\}$, \ $f(6) = \{a, b\}$\\[1mm]
            $f(2 \wedge 3) = f(1) = \varnothing = \{a\} \cap \{b\} = f(2) \cap f(3)$\\[1mm]
            $f(2 \vee 3) = f(6) = \{a, b\} = \{a\} \cup \{b\} = f(2) \cup f(3)$};
    \end{tikzpicture}
    \end{center}
    Nella parte $(\leftarrow)$ abbiamo usato solo $\wedge$: una biezione che rispetta $\wedge$ rispetta l'ordine, e quindi, per la parte $(\rightarrow)$, rispetta
    automaticamente anche $\vee$.
\end{spiegazione}

% ──────────────────────────────────────────────────────────────
\section{Il reticolo trirettangolo e il reticolo pentagonale}\label{sec:m3-n5}

Due piccoli reticoli di cinque elementi giocano un ruolo speciale: vedremo che sono proprio loro a impedire la distributività.

\dfn[m3-n5]{Reticolo trirettangolo e reticolo pentagonale}{
    Il \textbf{reticolo trirettangolo} e il \textbf{reticolo pentagonale} sono gli insiemi ordinati $\{a, b, c, d, e\}$ con i diagrammi di Hasse seguenti.
    \begin{center}
    \begin{tikzpicture}[hasse]
        \trirettangolo
        \node[font=\small] at (0,-0.65) {trirettangolo};
        \begin{scope}[xshift=6cm]
            \pentagonale
            \node[font=\small] at (0,-0.65) {pentagonale};
        \end{scope}
    \end{tikzpicture}
    \end{center}
    Nel trirettangolo $d$ è il minimo, $b$ il massimo e $a, c, e$ sono a due a due non confrontabili. Nel pentagonale $e$ è il minimo, $b$ il massimo,
    $e < d < c < b$ ed $e < a < b$, e $a$ non è confrontabile né con $c$ né con $d$.
}

\begin{concetto}
\begin{Proposizione}{Trirettangolo e pentagonale sono reticoli}{m3-n5-reticoli}
    \begin{prerequisiti}
        \dip{Definizione}{def:m3-n5}
        \dip{Definizione}{def:reticolo}
        \dip{Proposizione}{prop:minimo-inf}
    \end{prerequisiti}
    Il reticolo trirettangolo e il reticolo pentagonale sono reticoli limitati.
\end{Proposizione}
\begin{dimostrazione}
    Per due elementi confrontabili, inf e sup sono il minore e il maggiore (Proposizione~\ref{prop:minimo-inf}); restano le coppie non confrontabili.
    \begin{caso}{Trirettangolo}
        Le coppie non confrontabili sono formate da due elementi distinti tra $a$, $c$, $e$. Per ognuna l'unico maggiorante è $b$ e l'unico minorante è $d$:
        il sup è $b$ e l'inf è $d$.
    \end{caso}
    \begin{caso}{Pentagonale}
        Le coppie non confrontabili sono $\{a, c\}$ e $\{a, d\}$. Per entrambe l'unico maggiorante è $b$ ($c$ sta sopra $d$, ma non sopra $a$) e l'unico
        minorante è $e$ ($d$ sta sotto $c$, ma non sotto $a$): il sup è $b$ e l'inf è $e$.
    \end{caso}
    In entrambi i casi il minimo e il massimo esistono, quindi i reticoli sono limitati.
\end{dimostrazione}
\end{concetto}

\begin{spiegazione}
    Ecco le tavole delle operazioni del pentagonale, ottenute leggendo il diagramma (in arancione i risultati delle coppie non confrontabili).
    Nel pentagonale $a \vee d$ «salta» $c$: $c$ sta sopra $d$ ma non sopra $a$, quindi il primo elemento sopra entrambi è $b$.
    \begin{center}
    \begin{tikzpicture}[hasse]
        \pentagonale[n5a/.style=hv evid]
        \node[anchor=west] at (1.9,1.4) {\small
            \renewcommand{\arraystretch}{1.1}
            \begin{tabular}{c|ccccc}
                $\wedge$ & $e$ & $d$ & $c$ & $a$ & $b$ \\ \hline
                $e$ & $e$ & $e$ & $e$ & $e$ & $e$ \\
                $d$ & $e$ & $d$ & $d$ & \cellcolor{orange!25}$e$ & $d$ \\
                $c$ & $e$ & $d$ & $c$ & \cellcolor{orange!25}$e$ & $c$ \\
                $a$ & $e$ & \cellcolor{orange!25}$e$ & \cellcolor{orange!25}$e$ & $a$ & $a$ \\
                $b$ & $e$ & $d$ & $c$ & $a$ & $b$ \\
            \end{tabular}\qquad
            \begin{tabular}{c|ccccc}
                $\vee$ & $e$ & $d$ & $c$ & $a$ & $b$ \\ \hline
                $e$ & $e$ & $d$ & $c$ & $a$ & $b$ \\
                $d$ & $d$ & $d$ & $c$ & \cellcolor{orange!25}$b$ & $b$ \\
                $c$ & $c$ & $c$ & $c$ & \cellcolor{orange!25}$b$ & $b$ \\
                $a$ & $a$ & \cellcolor{orange!25}$b$ & \cellcolor{orange!25}$b$ & $a$ & $b$ \\
                $b$ & $b$ & $b$ & $b$ & $b$ & $b$ \\
            \end{tabular}};
    \end{tikzpicture}
    \end{center}
    Nel trirettangolo le tavole sono ancora più semplici: per due elementi distinti tra $a$, $c$, $e$ l'inf è sempre $d$ e il sup è sempre $b$.
\end{spiegazione}

% ──────────────────────────────────────────────────────────────
\section{Sottoreticoli}\label{sec:sottoreticoli}

Visto che i reticoli sono strutture algebriche, ha senso parlare delle loro sottostrutture.

\dfn[sottoreticolo]{Sottoreticolo}{
    Sia $(s, \rho)$ un reticolo. Una parte non vuota $t$ di $s$ si dice \textbf{sottoreticolo} di $(s, \rho)$ se è chiusa rispetto alle due operazioni
    reticolari: per ogni $x, y \in t$, $x \wedge y \in t$ e $x \vee y \in t$.
}

\begin{concetto}
\begin{Proposizione}{Un sottoreticolo è un reticolo}{sottoreticolo-reticolo}
    \begin{prerequisiti}
        \dip{Definizione}{def:sottoreticolo}
        \dip{Proposizione}{prop:ordine-indotto}
        \dip{Definizione}{def:inf-sup}
    \end{prerequisiti}
    Se $t$ è un sottoreticolo di $(s, \rho)$, allora $t$ con l'ordine indotto è un reticolo, e le sue operazioni sono le restrizioni a $t$ di $\wedge$ e $\vee$.
\end{Proposizione}
\begin{dimostrazione}
    Siano $x, y \in t$. L'elemento $x \wedge y$ sta in $t$ ed è un minorante di $\{x, y\}$; ogni minorante di $\{x, y\}$ in $t$ è anche un minorante in $s$, quindi
    sta sotto $x \wedge y$. Dunque $x \wedge y$ è l'estremo inferiore di $\{x, y\}$ anche in $t$. Dualmente $x \vee y$ è l'estremo superiore di $\{x, y\}$ in $t$.
\end{dimostrazione}
\end{concetto}

\begin{spiegazione}
    Il viceversa è falso: una parte può essere un reticolo con l'ordine indotto senza essere un sottoreticolo, perché i suoi inf e sup
    possono essere \textbf{diversi} da quelli di $s$. Nel reticolo $s = \{a, b, c, d, e\}$ qui sotto:
    \begin{center}
    \begin{tikzpicture}[hasse]
        \node[hv] (e) at (0,0) {$e$}; \node[hv evid2] (d) at (0,0.9) {$d$};
        \node[hv] (c) at (-0.9,1.8) {$c$}; \node[hv] (a) at (0.9,1.8) {$a$}; \node[hv] (b) at (0,2.7) {$b$};
        \draw[hlato] (e) -- (d) -- (c) -- (b); \draw[hlato] (d) -- (a) -- (b);
        \node[font=\small, align=center] at (0,-0.75) {il reticolo $s$:\\ $a \wedge c = d$};
        \begin{scope}[xshift=4.6cm]
            \node[hv evid] (e) at (0,0) {$e$}; \node[hv grigio] (d) at (0,0.9) {$d$};
            \node[hv evid] (c) at (-0.9,1.8) {$c$}; \node[hv evid] (a) at (0.9,1.8) {$a$}; \node[hv evid] (b) at (0,2.7) {$b$};
            \draw[hlato grigio] (e) -- (d) -- (c); \draw[hlato grigio] (d) -- (a);
            \draw[hlato] (c) -- (b) -- (a);
            \draw[hlato evid] (e) to[bend left=35] (c); \draw[hlato evid] (e) to[bend right=35] (a);
            \node[font=\small, align=center] at (0,-0.75) {$t = \{a, b, c, e\}$: in $t$ vale\\ $\inf\{a, c\} = e \neq d$};
        \end{scope}
        \begin{scope}[xshift=9.2cm]
            \node[hv grigio] (e) at (0,0) {$e$}; \node[hv evid2] (d) at (0,0.9) {$d$};
            \node[hv evid2] (c) at (-0.9,1.8) {$c$}; \node[hv evid2] (a) at (0.9,1.8) {$a$}; \node[hv evid2] (b) at (0,2.7) {$b$};
            \draw[hlato grigio] (e) -- (d);
            \draw[hlato] (d) -- (c) -- (b); \draw[hlato] (d) -- (a) -- (b);
            \node[font=\small, align=center] at (0,-0.75) {$\{a, b, c, d\}$:\\ è un sottoreticolo};
        \end{scope}
    \end{tikzpicture}
    \end{center}
    $t = \{a, b, c, e\}$ con l'ordine indotto è un reticolo (un «quadrato»), ma non è un sottoreticolo: $a \wedge c = d$ in $s$, e $d \notin t$.
    Invece $\{a, b, c, d\}$ è chiuso rispetto a $\wedge$ e $\vee$: è un sottoreticolo.
\end{spiegazione}

\begin{concetto}
\begin{Proposizione}{Gli intervalli chiusi sono sottoreticoli}{intervalli-sottoreticoli}
    \begin{prerequisiti}
        \dip{Definizione}{def:sottoreticolo}
        \dip{Definizione}{def:minoranti}
    \end{prerequisiti}
    Sia $(s, \rho)$ un reticolo e siano $x, y \in s$ con $x \rho y$. L'\textbf{intervallo chiuso} $[x, y] := \{z \in s \mid x \rho z \text{ e } z \rho y\}$ è un sottoreticolo.
\end{Proposizione}
\begin{dimostrazione}
    $[x, y] \neq \varnothing$ perché contiene $x$. Siano $z, w \in [x, y]$. $x$ è un minorante di $\{z, w\}$, quindi $x \rho (z \wedge w)$; inoltre $(z \wedge w) \rho z \rho y$.
    Dunque $z \wedge w \in [x, y]$. Dualmente $y$ è un maggiorante di $\{z, w\}$, quindi $(z \vee w) \rho y$, e $x \rho z \rho (z \vee w)$: anche $z \vee w \in [x, y]$.
\end{dimostrazione}
\end{concetto}

\begin{spiegazione}
    Un intervallo è il «pezzo di diagramma» compreso tra $x$ e $y$: inf e sup di due suoi elementi non possono uscire, perché $x$ resta sotto e $y$
    resta sopra. Nel cubo delle parti di $\{1, 2, 3\}$, l'intervallo $[\{1\}, \{1, 2, 3\}]$ (gli insiemi che contengono $1$) è un quadrato chiuso rispetto
    a $\cap$ e $\cup$:
    \begin{center}
    \begin{tikzpicture}[hasse]
        \cubotre[k1/.style=hv evid, k12/.style=hv evid, k13/.style=hv evid, k123/.style=hv evid]
        \draw[hlato evid] (k1) -- (k12) -- (k123); \draw[hlato evid] (k1) -- (k13) -- (k123);
        \node[font=\small, align=left, anchor=west] at (2.6,1.65) {ad esempio $\{1, 2\} \cap \{1, 3\} = \{1\}$ e $\{1, 2\} \cup \{1, 3\} = \{1, 2, 3\}$\\ restano nell'intervallo.
            Se invece non vale $x \rho y$, l'insieme\\ $\{z \mid x \rho z \text{ e } z \rho y\}$ è vuoto (per la transitività sarebbe $x \rho y$).};
    \end{tikzpicture}
    \end{center}
\end{spiegazione}

% ──────────────────────────────────────────────────────────────
\section{Complementi}\label{sec:complementi}

\dfn[complemento]{Elemento complementato, reticolo complementato}{
    Sia $(s, \rho)$ un reticolo limitato. Un elemento $x \in s$ si dice \textbf{complementato} se esiste $y \in s$ tale che
    \[ x \wedge y = \min(s) \qquad\text{e}\qquad x \vee y = \max(s); \]
    $y$ si dice un \textbf{complemento} di $x$ (e $x$ è un complemento di $y$). Un reticolo limitato si dice \textbf{complementato} se ogni suo elemento ha
    almeno un complemento.
}

\begin{esempio}{Complementi}
    \begin{itemize}
        \item $(\mbN, \leq)$ non è complementato: non è nemmeno limitato.
        \item In $(\mbN, \mid)$ (minimo $1$, massimo $0$) solo $0$ e $1$ sono complementati. Infatti $x \vee y = 0$ vuol dire $\mathrm{mcm}(x, y) = 0$, cioè $x = 0$ oppure
              $y = 0$; se $y = 0$, allora $x \wedge y = \mathrm{MCD}(x, 0) = x$ deve essere $1$. Quindi le uniche coppie di complementi sono $0$ e $1$.
        \item I divisori di $15$, cioè $(\{1, 3, 5, 15\}, \mid)$, formano un reticolo complementato: il complemento di $1$ è $15$ e quello di $3$ è $5$.
        \item $(\mcP(s), \subseteq)$ è complementato: un complemento di $x$ è $s \setminus x$, perché $x \cap (s \setminus x) = \varnothing$ e $x \cup (s \setminus x) = s$.
        \item Trirettangolo e pentagonale sono complementati, ma i complementi possono essere più d'uno: nel trirettangolo $a$ ha come complementi
              sia $c$ sia $e$; nel pentagonale $a$ ha come complementi sia $c$ sia $d$, mentre $c$ e $d$ hanno come unico complemento $a$.
    \end{itemize}
\end{esempio}

\begin{concetto}
\begin{Proposizione}{Complementi confrontabili}{complemento-confrontabile}
    \begin{prerequisiti}
        \dip{Definizione}{def:complemento}
        \dip{Proposizione}{prop:ordine-operazioni}
    \end{prerequisiti}
    Sia $(s, \rho)$ un reticolo limitato e sia $y$ un complemento di $x$. Se $x$ e $y$ sono confrontabili, allora $\{x, y\} = \{\min(s), \max(s)\}$. Quindi gli unici
    elementi confrontabili con un proprio complemento sono il minimo e il massimo.
\end{Proposizione}
\begin{dimostrazione}
    Se $x \rho y$, allora $x \wedge y = x$ e $x \vee y = y$ (Proposizione~\ref{prop:ordine-operazioni}); ma $x \wedge y = \min(s)$ e $x \vee y = \max(s)$, quindi $x = \min(s)$
    e $y = \max(s)$. Se $y \rho x$, scambiando i ruoli si ottiene $y = \min(s)$ e $x = \max(s)$. Viceversa, il minimo e il massimo sono confrontabili e sono
    complementi l'uno dell'altro.
\end{dimostrazione}
\end{concetto}

\begin{spiegazione}
    Un elemento e il suo complemento «si incontrano» solo nel minimo e «si raggiungono» solo nel massimo. Se uno dei due stesse
    sopra l'altro, il loro inf sarebbe quello più in basso e il loro sup quello più in alto: allora sarebbero proprio il minimo e il massimo.
    \begin{center}
    \begin{tikzpicture}[hasse]
        \cubotre[k1/.style=hv evid, k23/.style=hv evid, k0/.style=hv evid2, k123/.style=hv evid2]
        \node[font=\small, align=left, anchor=west] at (2.6,1.65) {in $\mcP(\{1, 2, 3\})$ il complemento di $\{1\}$ è $\{2, 3\}$ (in arancione):\\
            $\{1\} \cap \{2, 3\} = \varnothing$ e $\{1\} \cup \{2, 3\} = \{1, 2, 3\}$ (in verde).\\
            $\{1\}$ e $\{2, 3\}$ non sono confrontabili.};
    \end{tikzpicture}
    \end{center}
\end{spiegazione}

% ──────────────────────────────────────────────────────────────
\section{Reticoli distributivi}\label{sec:distributivi}

\dfn[distributivo]{Reticolo distributivo}{
    Un reticolo $(s, \rho)$ si dice \textbf{distributivo} se per ogni $a, b, c \in s$ valgono
    \[ \text{(D1)}\quad a \wedge (b \vee c) = (a \wedge b) \vee (a \wedge c) \qquad\qquad\qquad \text{(D2)}\quad a \vee (b \wedge c) = (a \vee b) \wedge (a \vee c). \]
}

\begin{concetto}
\begin{Proposizione}{Le due distributività sono equivalenti}{distributive-equivalenti}
    \begin{prerequisiti}
        \dip{Definizione}{def:distributivo}
        \dip{Teorema}{teo:proprieta-reticolari}
        \dip{Teorema}{teo:dualita-reticoli}
    \end{prerequisiti}
    In un reticolo, (D1) vale per ogni $a, b, c$ se e solo se (D2) vale per ogni $a, b, c$. Per sapere se un reticolo è distributivo basta quindi
    verificarne una.
\end{Proposizione}
\begin{dimostrazione}
    \facoltativa
    Supponiamo che valga (D1). Per ogni $a, b, c$:
    \begin{align*}
        (a \vee b) \wedge (a \vee c) &= \big((a \vee b) \wedge a\big) \vee \big((a \vee b) \wedge c\big) && \text{(D1)} \\
        &= a \vee \big((a \wedge c) \vee (b \wedge c)\big) && \text{assorbimento, commutatività e (D1)} \\
        &= \big(a \vee (a \wedge c)\big) \vee (b \wedge c) = a \vee (b \wedge c) && \text{associatività e assorbimento.}
    \end{align*}
    Quindi l'enunciato «se vale (D1), vale (D2)» è vero per ogni reticolo. Il suo duale è «se vale (D2), vale (D1)», che allora è vero per ogni
    reticolo per il principio di dualità.
\end{dimostrazione}
\end{concetto}

\begin{spiegazione}
    (D1) e (D2) sono una la duale dell'altra: ecco perché, dimostrata un'implicazione, l'altra viene gratis. Una verifica in $(\mbN, \mid)$
    con $a = 4$, $b = 6$, $c = 10$:
    \begin{passi}
        \item (D1): $4 \wedge (6 \vee 10) = \mathrm{MCD}(4, 30) = 2$ e $(4 \wedge 6) \vee (4 \wedge 10) = \mathrm{mcm}(2, 2) = 2$;
        \item (D2): $4 \vee (6 \wedge 10) = \mathrm{mcm}(4, 2) = 4$ e $(4 \vee 6) \wedge (4 \vee 10) = \mathrm{MCD}(12, 20) = 4$.
    \end{passi}
\end{spiegazione}

\begin{esempio}{Reticoli distributivi}
    \begin{itemize}
        \item $(\mcP(s), \subseteq)$ è distributivo: (D1) e (D2) sono le distributività di $\cap$ e $\cup$ (Proposizione~\ref{prop:distributivita-ins}).
        \item Ogni insieme totalmente ordinato è un reticolo distributivo. Dati $a, b, c$, per la commutatività di $\vee$ possiamo supporre $b \rho c$. Allora
              $b \vee c = c$; inoltre $a \wedge b$ sta sotto $a$ e sotto $b \rho c$, quindi $(a \wedge b) \rho (a \wedge c)$. Perciò
              $(a \wedge b) \vee (a \wedge c) = a \wedge c = a \wedge (b \vee c)$.
        \item $(\mbN, \mid)$ è distributivo (non lo dimostriamo). L'idea: per i numeri positivi MCD e mcm si calcolano esponente per esponente con
              $\min$ e $\max$ (Proposizione~\ref{prop:divisori-scomposizione}), e su ogni esponente vale la distributività appena vista per $(\mbN, \leq)$.
    \end{itemize}
\end{esempio}

\begin{concetto}
\begin{Teorema}{Unicità del complemento}{complemento-unico}
    \begin{prerequisiti}
        \dip{Definizione}{def:distributivo}
        \dip{Definizione}{def:complemento}
        \dip{Proposizione}{prop:neutri-reticolo}
    \end{prerequisiti}
    In un reticolo distributivo limitato ogni elemento ha al più un complemento.
\end{Teorema}
\begin{dimostrazione}
    Siano $y$ e $z$ complementi di $x$. Usando che $\max(s)$ è neutro per $\wedge$ e $\min(s)$ è neutro per $\vee$ (Proposizione~\ref{prop:neutri-reticolo}):
    \[ y = y \wedge \max(s) = y \wedge (x \vee z) \overset{\text{(D1)}}{=} (y \wedge x) \vee (y \wedge z) = \min(s) \vee (y \wedge z) = y \wedge z. \]
    Scambiando i ruoli di $y$ e $z$ si ottiene $z = z \wedge y$. Quindi $y = y \wedge z = z \wedge y = z$.
\end{dimostrazione}
\end{concetto}

\begin{spiegazione}
    In un reticolo distributivo il complemento, se c'è, è unico, e lo indicheremo con $x'$: in $(\mcP(s), \subseteq)$ è $s \setminus x$. Il teorema serve
    soprattutto «al contrario»: se un elemento ha \textbf{due} complementi, il reticolo non è distributivo. Nel trirettangolo $a$ ha i due
    complementi $c$ ed $e$, quindi il trirettangolo non è distributivo:
    \begin{center}
    \begin{tikzpicture}[hasse]
        \trirettangolo[m3a/.style=hv evid, m3c/.style=hv evid2, m3e/.style=hv evid2]
        \node[font=\small, align=left, anchor=west] at (1.9,1.1) {$a \wedge c = d = \min$ e $a \vee c = b = \max$: $c$ è un complemento di $a$;\\
            $a \wedge e = d = \min$ e $a \vee e = b = \max$: anche $e$ è un complemento di $a$.};
    \end{tikzpicture}
    \end{center}
\end{spiegazione}

\begin{concetto}
\begin{Proposizione}{Sottoreticoli di reticoli distributivi}{sottoreticoli-distributivi}
    \begin{prerequisiti}
        \dip{Definizione}{def:distributivo}
        \dip{Proposizione}{prop:sottoreticolo-reticolo}
    \end{prerequisiti}
    Ogni sottoreticolo di un reticolo distributivo è distributivo.
\end{Proposizione}
\begin{dimostrazione}
    Le operazioni del sottoreticolo $t$ sono le restrizioni di quelle di $s$ (Proposizione~\ref{prop:sottoreticolo-reticolo}). (D1) vale per ogni $a, b, c \in s$, quindi
    in particolare per ogni $a, b, c \in t$.
\end{dimostrazione}
\end{concetto}

\begin{spiegazione}
    Un'uguaglianza che vale per tutti gli elementi di $s$ vale anche per quelli di una sua parte, purché le operazioni restino le stesse:
    per questo serve un sottoreticolo e non una parte qualunque. Ad esempio, per $n \neq 0$ l'insieme dei divisori positivi di $n$ è un sottoreticolo di
    $(\mbN, \mid)$: se $x \mid n$ e $y \mid n$, allora $\mathrm{MCD}(x, y) \mid n$ e, poiché $n$ è un multiplo comune, $\mathrm{mcm}(x, y) \mid n$. Quindi i divisori di $12$,
    di $30$, di $210$ e così via formano reticoli distributivi.
\end{spiegazione}

\begin{concetto}
\begin{Proposizione}{Trirettangolo e pentagonale non sono distributivi}{m3-n5-non-distributivi}
    \begin{prerequisiti}
        \dip{Definizione}{def:distributivo}
        \dip{Proposizione}{prop:m3-n5-reticoli}
    \end{prerequisiti}
    Né il reticolo trirettangolo né il reticolo pentagonale sono distributivi.
\end{Proposizione}
\begin{dimostrazione}
    Basta trovare una terna per cui (D1) è falsa.
    \begin{caso}{Trirettangolo}
        $a \wedge (c \vee e) = a \wedge b = a$, mentre $(a \wedge c) \vee (a \wedge e) = d \vee d = d \neq a$.
    \end{caso}
    \begin{caso}{Pentagonale}
        $c \wedge (a \vee d) = c \wedge b = c$, mentre $(c \wedge a) \vee (c \wedge d) = e \vee d = d \neq c$.
    \end{caso}
\end{dimostrazione}
\end{concetto}

\begin{spiegazione}
    Seguiamo il calcolo per il pentagonale sul diagramma:
    \begin{center}
    \begin{tikzpicture}[hasse]
        \pentagonale[n5c/.style=hv evid, n5a/.style=hv evid2, n5d/.style=hv evid2]
        \node[font=\small, align=left, anchor=west] at (1.8,1.4) {
            \begin{minipage}{12.5cm}
            \begin{passi}
                \item a sinistra: $a \vee d = b$ (il primo elemento sopra $a$ e $d$), poi $c \wedge b = c$;
                \item a destra: $c \wedge a = e$ e $c \wedge d = d$, poi $e \vee d = d$;
                \item i risultati sono $c$ e $d$: «distribuire» $c \wedge$ sui due termini fa perdere il pezzo tra $d$ e $c$.
            \end{passi}
            \end{minipage}};
    \end{tikzpicture}
    \end{center}
    In alternativa: entrambi i reticoli hanno un elemento con due complementi, che in un reticolo distributivo non può esistere (Teorema~\ref{teo:complemento-unico}).
\end{spiegazione}

Il prossimo teorema dice che trirettangolo e pentagonale sono gli \textbf{unici} ostacoli alla distributività.

\begin{concetto}
\begin{Teorema}{Criterio di distributività di Birkhoff}{birkhoff}
    \begin{prerequisiti}
        \dip{Definizione}{def:distributivo}
        \dip{Definizione}{def:sottoreticolo}
        \dip{Definizione}{def:m3-n5}
        \dip{Teorema}{teo:isomorfismi-reticoli}
    \end{prerequisiti}
    \senzadim
    Un reticolo è distributivo se e solo se non ha sottoreticoli isomorfi al reticolo trirettangolo o al reticolo pentagonale.
\end{Teorema}
\end{concetto}

\begin{spiegazione}
    Una direzione si capisce subito: se un reticolo ha un sottoreticolo trirettangolo o pentagonale, questo non è distributivo, e allora
    non lo è nemmeno il reticolo (Proposizione~\ref{prop:sottoreticoli-distributivi}). Attenzione a cercare \textbf{sottoreticoli}: $(\mcP(\{1, 2, 3\}), \subseteq)$ è distributivo, eppure
    contiene una parte che, con l'ordine indotto, è un pentagono:
    \begin{center}
    \begin{tikzpicture}[hasse]
        \cubotre[k0/.style=hv evid, k1/.style=hv evid, k12/.style=hv evid, k3/.style=hv evid, k123/.style=hv evid, k13/.style=hv evid3,
                 k2/.style=hv grigio, k23/.style=hv grigio]
        \draw[hlato evid] (k0) -- (k1) -- (k12) -- (k123); \draw[hlato evid] (k0) -- (k3);
        \draw[hlato evid] (k3) -- (k123);
        \node[font=\small, align=left, anchor=west] at (2.6,1.65) {$t = \{\varnothing, \{1\}, \{1, 2\}, \{3\}, \{1, 2, 3\}\}$ con $\subseteq$ è un pentagono,\\
            ma non è un sottoreticolo: $\{1\} \cup \{3\} = \{1, 3\} \notin t$.\\
            In $t$ il sup di $\{1\}$ e $\{3\}$ è $\{1, 2, 3\}$, in $\mcP(\{1, 2, 3\})$ è $\{1, 3\}$.};
    \end{tikzpicture}
    \end{center}
\end{spiegazione}

\begin{esempio}{I divisori di 30 e i divisori di 12}
    \label{es:divisori-30-12}
    \begin{center}
    \begin{tikzpicture}[hasse]
        \node[hv] (1) at (0,0) {$1$};
        \node[hv evid] (2) at (-1.2,1.1) {$2$}; \node[hv evid2] (3) at (0,1.1) {$3$}; \node[hv] (5) at (1.2,1.1) {$5$};
        \node[hv] (6) at (-1.2,2.2) {$6$}; \node[hv evid2] (10) at (0,2.2) {$10$}; \node[hv evid] (15) at (1.2,2.2) {$15$};
        \node[hv] (30) at (0,3.3) {$30$};
        \draw[hlato] (1) -- (2); \draw[hlato] (1) -- (3); \draw[hlato] (1) -- (5);
        \draw[hlato] (2) -- (6); \draw[hlato] (2) -- (10); \draw[hlato] (3) -- (6); \draw[hlato] (3) -- (15); \draw[hlato] (5) -- (10); \draw[hlato] (5) -- (15);
        \draw[hlato] (6) -- (30); \draw[hlato] (10) -- (30); \draw[hlato] (15) -- (30);
        \node[font=\small, align=center] at (0,-0.85) {divisori di $30$: distributivo e complementato,\\ il complemento di $x$ è $30/x$};
        \begin{scope}[xshift=8.4cm]
            \node[hv] (1) at (0,0) {$1$}; \node[hv evid3] (2) at (-0.8,1.1) {$2$}; \node[hv evid2] (3) at (0.8,1.1) {$3$};
            \node[hv evid2] (4) at (-1.6,2.2) {$4$}; \node[hv evid3] (6) at (0,2.2) {$6$}; \node[hv] (12) at (-0.8,3.3) {$12$};
            \draw[hlato] (1) -- (2) -- (4) -- (12); \draw[hlato] (1) -- (3) -- (6) -- (12); \draw[hlato] (2) -- (6);
            \node[font=\small, align=center] at (-0.4,-0.85) {divisori di $12$: distributivo, ma non complementato:\\ $2$ e $6$ (in rosso) non hanno complemento};
        \end{scope}
    \end{tikzpicture}
    \end{center}
    Entrambi sono distributivi, perché sottoreticoli di $(\mbN, \mid)$. Nei divisori di $30$ ogni elemento ha complemento: ad esempio $2' = 15$ (in arancione)
    e $3' = 10$ (in verde). Nei divisori di $12$ invece $1' = 12$ e $3' = 4$ (in verde), ma $2$ non ha complemento: $2 \wedge y = 1$ vuol dire $y$ dispari, cioè $y \in \{1, 3\}$, e
    $2 \vee 1 = 2$, $2 \vee 3 = 6$, mai $12$. Lo stesso vale per $6$, il cui unico $y$ con $6 \wedge y = 1$ è $y = 1$. Un reticolo distributivo e complementato come
    i divisori di $30$ si chiama \textbf{booleano}: è il protagonista del prossimo capitolo.
\end{esempio}
